\documentclass{article}
\usepackage{tikz}
\usetikzlibrary{angles, quotes}
\usepackage{amssymb}
\usepackage{pgfplots}
\pgfplotsset{width=10cm,compat=1.9}
\usepgfplotslibrary{statistics}
\usepackage[left=2.54cm,right=2.54cm,top=2.54cm,bottom=2.54cm]{geometry}
\usepackage{tasks} %Paquete necesario para la producción de items horizontales para algunos ejercicios de tarea empleando el comando /begin{multicols}}{#}.../end{multicols} para ayudarnos a reducir las páginas
\usepackage{pdflscape} %comando que permite cambiar a orientación horizontal las hojas%
\usepackage{soul} %Paquete para permitir la compilación de acentos usando el comando {\hl{}}}\\
\sethlcolor{yellow(munsell)} %comando para definir el color para el subrayado usando el comando \hl
\usepackage[utf8]{inputenc}
\usepackage[spanish]{babel}
\usepackage{icomma}
\usepackage{ragged2e}
\usepackage{siunitx}
\usepackage{url}
\usepackage[colorlinks=true, urlcolor=blue,  linkcolor=black, citecolor=green]{hyperref}
\usepackage{pdfpages}
\usepackage{blindtext} %Paquete que genera texto ficticio (dummy text) con el comando: \blindtext
\usepackage{cancel} %in the preamble gives you four different modes of striking through: \cancel{text to cancel} draws a diagonal line (slash) through its argument, \bcancel{text to cancel} uses the negative slope (a backslash), \xcancel{text to cancel} draws an X (actually \cancel plus \bcancel), \cancelto{〈value〉}{〈expression〉} draws a diagonal arrow through the 〈expression〉pointing to the 〈value〉 (math-mode only)
\usepackage{csquotes}
\usepackage{afterpage}
\usepackage{parskip} 
\usepackage{float}
\usepackage{enumitem}
\usepackage{multicol}%Paquete que permite la creación aislada de columnas de texto. De esta forma se reduce la cantidad de páginas en nuestro documento
\usepackage{lipsum} %Paquete que genera texto de relleno al igual que \usepackage{blindtext}, se usa con el comando \lipsum[#-#] los #(hashtags) delimitan el número de páginas que deseamos ocupar del paquete lipsum para usarlos como relleno. 

% Definir el comando personalizado para la nota
\newcommand{\mynote}[1]{%
\begin{tikzpicture}[baseline=-0.75ex]
\node[draw, circle, fill=Apple Green!60, inner sep=2pt] (note) {\textbf{Nota:}};
\end{tikzpicture}%
\ \textit{#1}%
}

\newenvironment{Figura}
{\par\medskip\noindent\minipage{\linewidth}}
{\endminipage\par\medskip}
\usepackage{caption}
\usepackage[
backend=biber,
sorting=none,
url=true
]{biblatex} %bibliografía
\addbibresource{biblio.bib}
% Establecer el espacio entre las entradas de la bibliografía
\setlength{\bibitemsep}{1\baselineskip} % Puedes ajustar el valor según tus preferencias
\usepackage{amsmath,amsthm,amssymb,amsfonts}
\usepackage{pifont} %Permite la compilación del comando \xmark: es el símbolo de la tachita
\usepackage{mathtools} %permite la compilación de símbolos para matrices
\usepackage{empheq} %Paquete que se relaciona con \usepackage[most]{tcolorbox} para la creación de cuadros/cajas de colores para delimitar los resultados matemáticos o líneas de texto usando la línea de comando: \begin{empheq}[box={\mymath[colback=orange(webcolor)!70,drop lifted shadow]}]{equation*} \end{empheq}
\usepackage[most]{tcolorbox} %Paquete que se relaciona con \usepackage{empheq} para la creación de cuadros/cajas de colores para delimitar los resultados matemáticos o líneas de texto
\renewcommand{\qedsymbol}{$\blacksquare$}
\usetikzlibrary{positioning,decorations.pathreplacing} %Paquete que permite la compilación de llaves para cuadros sinópticos (brace diagram) 
\usepackage{schemata} %The schemata package is designed just for these brace diagrams
\usepackage{booktabs}

\newcommand\AB[2]{\schema{\schemabox{#1}}{\schemabox{#2}}} %comando o paquete necesario para crear cuadros sinópticos

\newtcbox{\mymath}[1][]{%
nobeforeafter, math upper, tcbox raise base,
enhanced, colframe=blue!30!black,
colback=blue!30, boxrule=1pt,
#1} %Comando importante para encasillar los resultados matemáticos en cajas de diferentes colores, se relaciona con los paquetes \usepackage{empheq} y \usepackage[most]{tcolorbox} 

\newenvironment{sysmatrix}[1]
{\left(\begin{array}{@{}#1@{}}}
{\end{array}\right)}
\newcommand{\ro}[1]{%
\xrightarrow{\mathmakebox[\rowidth]{#1}}%
}
\newlength{\rowidth}% row operation width
\AtBeginDocument{\setlength{\rowidth}{3em}}  %Comando importante para laproducción de lineas de operaciones entre matrices, método de Gauss o Gauss Jordan se relaciona con \newenvironment{sysmatrix}

%formato para cambiar el horario a español
\usepackage[spanish]{datetime2}
\DTMsetdatestyle{spanish}

\renewcommand{\today}{\DTMdisplaydate{\the\year}{\the\month}{\the\day }{-1}}
%%%%%%%%%%%%%%%%%%%%%%%%%%%%%%%%%%%%%%%%%

\usepackage{datetime}
\newcommand{\mycurrenttime}{\xxivtime}

%%%%%%%%%%%%%%%%%%%%%%%%%%%%%%%%%%%%%%%%%%%%%%%%%%%%%%%%%%%%%%%%%%%%%%%%%%%%%%%%%%%%%%%%%%%%%%%%%%%%%%%%%%%%%%%%%%%%%%%%%%%%%%%%%%%%%%%%%%%%%%%%%


%%%%%%%%%%%%%Caja de color para el título%%%%%%%%%%%%%%%%%%%%%%%%%%%%%%%%%%%%%%%%%%%%%%%%%%%%%%

\definecolor{myframecolor}{RGB}{1, 45, 75} % Prussian Blue
\definecolor{myboxcolor}{RGB}{0, 107, 92} % Apple Green

\newtcolorbox{mybox}{
enhanced,
colback=myboxcolor!7, % Color de fondo del cuadro
colframe=myframecolor, % Color del marco del cuadro
arc=0pt,
boxrule=1pt,
borderline west={2mm}{-10mm}{myframecolor}, % Borde en el lado izquierdo
sharp corners=southwest,
width=\linewidth,
before=\par\vspace{\bigskipamount}, % Espacio antes del cuadro
after=\par\vspace{\bigskipamount} % Espacio después del cuadro
}


\newcommand{\euler}{e} %Comando para producir letra e de euler
\newenvironment{remark}{\par\vfill\footnotesize % Vertical white space above the remark and smaller font size
\begin{list}{}{
	\leftmargin=80pt % Indentation on the left
	\rightmargin=60pt}\item\ignorespaces % Indentation on the right
\makebox[-2.5pt]{\begin{tikzpicture}[overlay]
		\node[draw=Horizon!60,line width=2.5pt,circle,fill=Horizon!25,font=\sffamily\bfseries,inner sep=4pt,outer sep=0pt] at (-19pt,5pt){\textcolor{Horizon}{Nota}};\end{tikzpicture}} % Blue Nota in a circle
\advance\baselineskip -1pt}{\end{list}\vskip5pt} % Tighter line spacing and white space after remark
\usepackage{graphicx}

\usepackage{titling}

\input{colores.tex}

%%Producir signos de cita textual``Asignatura''

\newcommand{\xmark}{\ding{55}} %%%comando que genera la tachita

\newenvironment{MyColorPar}[1]{%
\leavevmode\color{#1}\ignorespaces%
}{%
}%

%%%%%%%%%%%%%%%%%%% Título de la tarea, Nombre de alumno

\title{Clase \# 35 - Mecánica Estadística}
\author{\emph{JA} $\boldsymbol{\mid}$}
\date{\today}

%%%%%%%%%%%%%%%%%%% Datos de la Materia

\usepackage{fancyhdr}
\fancypagestyle{plain}{%  the preset of fancyhdr 
\fancyhf{} % clear all header and footer fields
\fancyfoot[R]{
\includegraphics[width=2cm]{LogoFCFMBUAP (1).png}}
\fancyfoot[L]{{\bfseries{\thedate{}}} a las {\bfseries{\mycurrenttime{}}} horas{} (GMT-6, H. Puebla de Zaragoza, Pue)}
\fancyhead[L]{Mecánica Estadística}
\fancyhead[R]{\theauthor}
}
\makeatletter
\def\@maketitle{%
\newpage
\null
\vskip 1em%
\begin{center}%
\let \footnote \thanks
{\LARGE \@title \par}%
\vskip 1em%
%{\large \@date}%
\end{center}%
\par
\vskip 1em}
\makeatother

\usepackage{lipsum}  
\usepackage{cmbright}


%%%%%%%%%%%%%%%%%%%%%%%%%%%%%%%%%%%%%%%%%%%%%%%%%%%%%%%%%%%%%%%%%%%%%%%%%%%%%%%%%%%%%%%%%%%%%%%%%%%%%%%%%%%%%%%%%%%%%%%%%%%%%%%%%%%%%%%%%%%%%%%%%%%%%%%%%%%%%%%%%%%%%%%%%%%%%%%%%%%%%%%%%%%%%%%%%%%%%%

\begin{document}



\maketitle


%%%%%%%%%%%%%%%%%%% Datos del profesor, campus y aula

\begin{mybox}
\noindent\begin{tabular}{@{}ll}
	{\bfseries{Profesor}} & Dr. Jorge Velázquez Castro\\
	\textcolor{prussianblue}{\bfseries{Campus}} / \textcolor{trueblue}{\bfseries{Facultad}}  & \textcolor{prussianblue}{\bfseries{C.U. BUAP}} /  \textcolor{trueblue}{\bfseries{FCFM}} \\
	\textcolor{prussianblue}{\bfseries{Edificio}} /  \textcolor{trueblue}{\bfseries{Salón}}    & 	\textcolor{prussianblue}{\bfseries{EMA4}} /  \textcolor{trueblue}{\bfseries{407}}   
\end{tabular} 
\end{mybox} \vspace{1cm}

\tableofcontents

\vspace{2cm}


\section{Recordando}

\[V_{N}(r) = \frac{r^{N}\pi^{\frac{N}{2}}}{\frac{N}{2}\Gamma \left(\frac{N}{2}\right)}, \quad \text{con} \quad C_{N}= \frac{\pi^{\frac{N}{2}}}{\frac{N}{2}\Gamma \left( \frac{N}{2} \right)}  \]


\[ S_{N}(r) = C_{N}Nr^{N-1}  \]

\[ 2mE=\sum\limits_{i=1}^{3N} p_{i}^{2}  \]

\textbf{Ecuación de la esfera en 3N dim de radio}

\[ \sqrt{2mE} \]


La superficie de la esfea $E=H$ es:


\[ S_{3N}\left( \sqrt{2mE} \right)   \]



Si fijamos el valor de una dimensión $p_{i}=p$ al valor $p$ tenemos:


\[ 2mE-p^{2} = \sum\limits_{i=1}^{3N-1}p_{1}^{2}   \]

Por lo tanto su superficie es proporcional a:


\[ S_{3N-1} \left(  \sqrt{ 2mE-p^{2}  }   \right)   \]

La probabilidad de encontrar al sistema o subsistema con un valor $p$ es:


\[ \rho_{C}(P) = \frac{ S_{3N-1} \left( \sqrt{ 2mE-p^{2} }  \right)dP  }{S_{3N} \left( \sqrt{2mE }  \right) } \]


Veamos:

\begin{equation}
\begin{split}
 S_{3N-1} \left(  \sqrt{2mE-p^{2} }  \right) &= C_{N}N\left( 2mE-p^{2} \right)^{ \frac{N-1}{2} }  \\ 
 &= \frac{ C_{N}N(2mE)^{\frac{N}{2}} }{(2mE-p^{2})^{\frac{1}{2}}} \left( 1 -  \frac{p^{2}}{2mE} \right)^{\frac{N}{2}} 
\end{split}
\end{equation}



En el límite termodinámico:

\[ N \to \infty; E \to \infty \quad \text{de modo que} \quad \to \frac{E}{N} \to e, \quad \text{finito}\]


Tenemos:


\[ \rho_{c}(p) \longrightarrow C_{0} \left(1- \frac{p^{2}}{2m\underbrace{\frac{E}{N}}_{e}N} \right)  \]

\[\to C_{0}\left( e^{\frac{-p^{2}}{2me}}\right)^{\frac{3}{2}}\]


Tenemos:



\[ \rho_{c}(p) \underbrace{\overbrace{\longrightarrow}^{N \to \infty}}_{E\to \infty} C_{0}e^{- \frac{3}{2}\frac{p^{2}}{2me}}  \]

Como:


\[ \boxed{e=\frac{3}{2} K_{B}T } \]

Entonces:


\[  \rho_{c}(p)  \underbrace{\overbrace{\longrightarrow}^{N \to \infty}}_{E\to \infty} C_{0}e^{ - \frac{1}{K_{B}T} \left( \frac{p^{2}}{2m}\right) } \]


\subsection{N-Osciladores cuánticos independientes}

(Cristal Einstein)


Tenemos:

\[ \epsilon_{n} = \hbar w \left( n+\frac{1}{2}   \right); \quad n=0,1,2     \]

Sabemos que:

\[ Z(T,V,N) = \left[ Z(T,V,1)  \right]^{N} \]

Además:


\[ Z(T,V,1) = \sum\limits_{n=0}^{\infty} e^{-\beta\epsilon_{n}}   \]


Al sustituir la expresión $\epsilon_{n}$

Tenemos:

\begin{align}
Z(T,V,1) &= \sum_{n=0}^{\infty} e^{-\beta \hbar \omega\left(n+\tfrac{1}{2}\right)} \\
&= e^{-\beta \hbar \omega/2}\sum_{n=0}^{\infty} \left(e^{-\beta \hbar \omega}\right)^n \\
&= \frac{e^{-\beta \hbar \omega/2}}{1 - e^{-\beta \hbar \omega}}
= \frac{1}{2\sinh\!\left(\tfrac{\beta \hbar \omega}{2}\right)}
\end{align}


% Partición de un oscilador
\[
Z(T,V,1)=\frac{e^{-\beta \hbar \omega/2}}{1-e^{-\beta \hbar \omega}}
\]

% Energía libre para N osciladores no interactuantes
\[
F(T,V,N) = -N k_B T \ln Z(T,V,1)
= N\frac{\hbar \omega}{2} + N k_B T \ln\!\big(1 - e^{-\beta \hbar \omega}\big)
\]



o bien como:


\[  Z(T,V,1) = \frac{1}{ e^{ \frac{\beta \hbar w}{2}}  - e^{ \frac{-\beta \hbar w}{2}}  } = \frac{1}{  2\sen{h}\left(  \frac{\beta \hbar w}{2}  \right)  } \]

Tenemos:


\[F(T,V,N) = NK_{B}T \ln \left[ 2\sen{h} \left( \frac{\beta \hbar w}{2}  \right)   \right]  \]


Las ecuaciones de estado:

\[S = - \frac{\partial F}{\partial T} = NK_{B}T \left[  \frac{\beta \hbar w}{2} \cot{h} \left( \frac{\beta w \hbar}{2}  \right)  - \ln \left[   \right] \right]   \]


o bien:

\[ S = - \frac{\partial F}{\partial T} = NK_{B} \left[  \frac{ \beta \hbar w }{e^{\beta \hbar w} -1 } - \ln \left[ 1  - e^{-\beta \hbar w}\right]\right] \]

\[ \mu = \frac{\partial F}{\partial N} \]


También:

\[ \langle E \rangle = U = - \frac{\partial}{\partial \beta} \ln{z} \]

\[ U = N\hbar w \left[  \frac{1}{2} + \frac{1}{e^{\beta \hbar w} -1}  \right] \]

Tenemos:

\[\langle\epsilon_{n} \rangle = \langle \frac{E}{N}  \rangle = \frac{ \langle E \rangle }{N} \]

Tenemos:



\[ \langle\epsilon_{n} \rangle = \hbar w \left[  \frac{1}{2} + \frac{1}{ e^{\beta\hbar w}-1}  \right]  \]


POr otro lado:

\[ \langle \epsilon_{n} \rangle = \hbar w \left(  \frac{1}{2} + \langle n \rangle \right) \]

Tenemos:


\[ \boxed{ \langle n \rangle  = \frac{1}{e^{\beta \hbar w}-1} } \]



\begin{figure}[H]
    \centering
	    \includegraphics[width=0.7\textwidth]{1.png}
    \caption{.}   
    \label{fig:hamiltoniana_transformacion}
\end{figure}

\[ \beta \hbar w = \frac{\hbar w}{K_{B}T} \]


Recordemos:


\[ U = n\bar w \left[  \frac{1}{2} + \frac{1}{ e^{\beta \hbar w} -1}  \right] \]

Entonces:


\[ U \approx \frac{N}{2}\hbar w + N\hbar w \frac{1}{ 1+\beta \hbar w + \frac{1}{2} \left( \right)^{2}   } \]


Donde se usó:


\[ \frac{1}{1+x} \approx 1-x+x^{2}  \]

Por lo tanto:

\[ C_{v} = NK_{B}  \]

Para altas temperaturas $T>> \Theta_{E}$


El cálculo exacto (cuántico)

\[ C_{v} = NK_{B}T\left( \beta \hbar w \right)^{2} \frac{e^{\beta \hbar w}}{( e^{\beta \hbar w}-1 )^{2}} \]

\begin{figure}[H]
    \centering
	    \includegraphics[width=0.7\textwidth]{2.png}
    \caption{.}   
    \label{fig:hamiltoniana_transformacion}
\end{figure}

Se grafica

 \[  \frac{U}{N\hbar w } vs \frac{KT}{\hbar w} \]
 
 Siguiente clase vemos el cálculo del clásico.


\end{document}
